\section{参数未知，而统计量随机}
\label{sec:11-Mathematical-Statistics/01-Statistical-Models/02}

一次模型评测：测试集 2000 条，准确率 \(0.912\)，写进周报。两周后有人按同样的流程重新采了 2000 条，跑出 \(0.897\)。会上的问题是模型有没有退化。

上一节的语言可以立刻回答一半：采集流程没变，那两批数据来自同一个 \(P_\theta\)，模型的权重也一个比特都没动。变的只有抽到的那 2000 条。\(0.912\) 和 \(0.897\) 不是模型的两个属性，是同一个属性在两次抽样下留下的两个影子。

麻烦在于这两个数字在报表里长得跟模型的其他属性一模一样——都是一个小数，都带三位精度，都可以画进折线图看趋势。要把它们区分开，需要先承认报表里并排放着两类完全不同的对象：一类固定不动而你看不见，另一类你看得一清二楚而它一直在动。

\begin{center}
\begin{tikzpicture}[x=1.7cm,y=0.62cm,font=\small,>=stealth]
  \draw[->,black!55] (-2.35,0) -- (2.45,0) node[right,font=\footnotesize] {\(\hat\theta\)};
  \draw[dashed,black!70,thick] (0,-0.2) -- (0,5.3);
  \node[font=\footnotesize,black!75,anchor=north] at (0,-0.35) {\(\theta\)（固定，看不见）};
  \foreach \px/\py in {-0.9/0.55, 0.4/0.95, -0.2/1.35, 1.3/1.75, -1.5/2.15, 0.7/2.55,
                       -0.6/3.35, 1.0/3.75, -0.35/4.15, 0.55/4.55, -1.1/4.95}
    \fill[blue!55!black] (\px,\py) circle (2.2pt);
  \fill[red!70!black] (0.45,2.95) circle (3.0pt);
  \node[font=\footnotesize,red!70!black,anchor=west] at (0.62,2.95) {手上这一次};
  \draw[<->,black!55] (-1.5,5.7) -- (1.3,5.7);
  \node[font=\footnotesize,black!70,anchor=south] at (-0.1,5.78) {这团点的宽度就是标准误};
\end{tikzpicture}

\emph{只有那条竖线不动；报表上的数字是这些点里的一个，而你不知道是哪一个。}
\end{center}

图里每一个蓝点都是一次完整重复：重新采一批数据，重新跑一遍评测，得到一个数。红点是实际发生的那一次。竖线的位置由生产环境的真实准确率决定，它不随抽样改变，也永远不会直接出现在任何一张表里。频率学派框架下的全部不确定性都来自这张图：\(\theta\) 从不摇摆，散开的是点。后面的偏差、方差、标准误、覆盖率，说的都是这团点相对那条竖线的几何关系。

\subsection{统计量必须能从数据算出来}

固定不动的那一类叫参数。在经典频率学派表述里，\(\theta\) 是一个固定但未知的常数，换一次样本它不会改变数值。会改变的是从样本里算出来的任何量，这些量叫统计量。

统计量是样本的函数，

\[
T=T(X_1,\ldots,X_n),
\]

并且定义式中不许出现未知参数。样本均值 \(\bar X=\frac1n\sum_i X_i\) 是统计量：\(n\) 已知，每个观测值可以读出来，代进去就有结果。\(\bar X-\theta\) 不是统计量，你写不出它的数值；\(\sigma\) 未知时 \(\bar X/\sigma\) 也不是。这条要求看着琐碎，却在后面反复出场——一个量能不能当场从数据算出来，决定了它能不能充当检验统计量、能不能出现在区间的两端。

约束是对未知参数的，不是对所有常数的。已知的样本量、试验设计矩阵、原假设指定的 \(\mu_0\)、事先固定的超参数，都可以进公式。\(\sigma\) 已知时

\[
\frac{\bar X-\mu_0}{\sigma/\sqrt n}
\]

在检验 \(H_0:\mu=\mu_0\) 时可以直接计算，因为 \(\mu_0\) 是假设指定的值，不是未知真值。

需要小心的是那些看着像常数、其实是从数据里估出来的量。用样本方差替换 \(\sigma^2\)、用交叉验证选出 \(\lambda\)、用同一批数据先挑变量再报系数——这些做法把一个随机量当成了已知常数带进第二步。第一步引入的波动被抹掉了，第二步的标准误因此偏小，区间偏窄。模型选择之后的推断、数据驱动阈值、经验 Bayes 都踩在这条线上，后果是名义覆盖率与实际覆盖率对不上。

统计量还顺带做了一件事：压缩。\(n\) 维样本被压成几个数，样本均值扔掉了顺序、极端值和离散程度，只留下位置。丢掉的东西里含不含参数信息，决定了这次压缩是无损的还是有损的。Bernoulli 模型里成功总数 \(S=\sum_i X_i\) 留住了关于 \(\theta\) 的全部似然信息，知道成功落在哪几个位置对推断 \(\theta\) 没有额外帮助；但只留下均值就估不出方差了，离散程度的信息不在均值里。这条线索的严格版本见\hyperref[sec:11-Mathematical-Statistics/03-Sufficient-Statistics-and-Information/01]{``充分统计量保留了哪些参数信息''}。

\subsection{估计量是带着任务的统计量}

统计量这个词很宽松：样本最大值是统计量，第一个观测 \(X_1\) 也是统计量。用 \(X_1\) 去估 \(\theta\) 完全合法，只是糟糕。当一个统计量被指派去估计某个参数或参数的函数时，我们叫它估计量，记作

\[
\hat\theta=T(X_1,\ldots,X_n).
\]

观测到具体数据之后它取到一个数，

\[
\hat\theta_{\mathrm{obs}}=T(x_1,\ldots,x_n),
\]

这个数叫估计值。同一个公式，观测前是随机变量，观测后是一个不再变化的数。

大小写这层区分立刻能挡住一句常见的错话。``\(0.912\) 这个估计值的方差是多少''问不通：\(0.912\) 是个定数，定数没有方差。有方差的是产生它的那条规则在重复抽样中的表现，也就是图里那团点的宽度。同理，偏差不是某次估计错了多少——那个偏离你既算不出来也看不见——而是这团点的中心与竖线之间的系统性错位。把这两句话记牢，\hyperref[sec:11-Mathematical-Statistics/02-Point-Estimation/01]{``偏差、方差与均方误差''}里的分解就只是算术了。

回到那 2000 条测试样本。设第 \(i\) 条判对与否为 \(X_i\overset{\iid}{\sim}\Bernoulli(\theta)\)，其中 \(\theta\) 是模型在整个数据生成分布上的准确率。判对的条数 \(S=\sum_i X_i\) 是统计量，服从 \(\Binomial(n,\theta)\)，重复评测时它会取到 1824、1794、1841 这些值；\(\hat\theta=S/n\) 是估计量，同样随机；这一次看到 \(S=1824\)、\(n=2000\)，于是估计值等于 \(0.912\)。四层对象里只有第一层不动，而报表上写的是第四层。

\subsection{一个数不足以描述这团点}

于是``该报什么''有了一个可以计算的答案。在上面的模型下，

\[
\begin{aligned}
\Var(\hat\theta)&=\frac{\theta(1-\theta)}{n},\\
\SE(\hat\theta)&=\sqrt{\frac{\theta(1-\theta)}{n}}\approx\sqrt{\frac{0.912\times0.088}{2000}}\approx0.0063.
\end{aligned}
\]

条件对一遍：第一行只用到独立与同分布，独立让方差可加，同分布让每项方差相同；第二行把未知的 \(\theta\) 换成了 \(\hat\theta\)，这一步是近似，在 \(n\) 不太小、\(\theta\) 不太贴近 0 或 1 时才站得住。若测试样本按用户成组，独立性失效，这个 \(0.0063\) 会明显偏小。

\(0.0063\) 就是图中那团点的宽度尺度。据此报告 \(0.912\pm0.012\)，或者写成区间 \([0.900,0.924]\)，比单报 \(0.912\) 多告诉读者一件事：这个位置本身有多少晃动余地。而两周后的 \(0.897\) 与 \(0.912\) 相差 \(0.015\)，两批数据独立时差值的标准误约为 \(0.0093\)，这个差距还不到两个标准误——把它读成模型退化，等于把图里两个蓝点之间的正常间距当成了竖线的移动。至于``不到两个标准误''为什么就该按住不表，那是下一节的事。

三位小数的准确率、四位小数的 AUC、精确到 \(0.1\%\) 的转化率，这些数字的显示精度和它们的可信精度经常差一个数量级。报点估计还是报区间，取决于读者会拿它做什么决定：排个内部优先级，点估计够用；据此上线或回滚，就得让波动一起出现在报表上。

\subsection{要估的那个量未必是模型里的 \(\theta\)}

参数不一定是一个数。正态模型的参数是位置与尺度这一对 \((\mu,\sigma^2)\)，回归的参数是系数向量 \(\beta\)，协方差结构的参数是矩阵 \(\Sigma\)，非参数问题里参数干脆是整条密度或整条生存曲线。要报告的量也常常是参数的一个变换 \(g(\theta)\)：两组均值之差、风险比、提升度。

多数分析只关心其中一小块，其余未知量称为干扰参数。估计正态均值时你并不想报告 \(\sigma^2\)，但 \(\sigma^2\) 会通过标准误进到你的区间里；上面那个 \(0.0063\) 也一样，它依赖于未知的 \(\theta\)。干扰参数不必报告，却必须被处理，因为把它当作已知会让最终结论的变异性被系统性低估。

还有一类任务从一开始就不在估参数。预测下一条样本会不会判对，和估计模型在总体上的准确率，是两个不同的问题：前者的误差来自单条样本的随机性，后者的误差来自这团点的散布，两者的量级差得很远。同一个模型里，最优的估计量未必是最优的预测规则，评价标准也不该混用。

\subsection{Bayesian 把随机性放在另一边}

图里那条竖线之所以不动，是频率学派的一个选择，不是自然律。Bayesian 的做法是给 \(\theta\) 指定先验分布，让参数本身成为随机变量；此时``看不见''的那一层也有了分布，观测数据之后更新为后验。

两套框架的概率陈述指向不同的对象。频率学派说的是估计程序在重复抽样中的行为——竖线不动，点在散；Bayesian 说的是给定这一份数据后参数的条件分布——点已经落定，分布画在竖线的位置上。两种说法各自自洽，混写在同一个句子里就会出乱子，比如把置信区间读成``参数以 \(95\%\) 的概率落在其中''。这句话在频率学派里是错的，在 Bayesian 里换个名字才是对的，具体见\hyperref[sec:11-Mathematical-Statistics/04-Interval-Estimation/01]{``置信度属于程序，而非这次参数''}。

现在还剩一个空缺。整节都在用``这团点的散布''说话，可我们只做了一次实验，图里除了红点，其余的蓝点一个也没有发生过。那些点的位置究竟从哪里来，凭什么可以拿它们算出 \(0.0063\)？下一节回答这个问题。
